\documentclass[12pt,a4paper]{article}

\usepackage[german]{babel}
\usepackage[utf8]{inputenc}
\usepackage[T1]{fontenc}
\usepackage{enumerate}
\usepackage[absolute]{textpos}
\usepackage[a4paper, top=2cm, left=2cm, bottom=2cm, right=2cm]{geometry}
\usepackage{parskip}

\usepackage{graphicx}

\begin{document}
\thispagestyle{empty}

\begin{textblock*}{\paperwidth}(1.6cm,0.48cm) %{width}(horizontal, vertical)
\includegraphics[scale=0.67]{Logo_Mathematik_CMYK.pdf}
\end{textblock*}

\vspace*{1.5cm}

\begin{center}
\Huge{\textbf{Svolgimento dell'esame}}\\
\large{OSM 2020}
\end{center}

\bigskip

Le informazioni sottostanti regolano lo svolgimento della prova preliminare per le OSM.

\section{Indicazioni sul materiale}
\begin{itemize}
\item Necessario per scrivere e per disegnare (compasso, riga, squadretta).
\item Usare eslcusivamente fogli assolutamente privi di qualsiasi scritta (anche se cancellata). 
\item Nessun supporto elettronico (computer, telefoni, calcolatrici,...).
\item Bevande e cibo sono ammessi (nel limite del buon senso e se questi non emettono odori sgradevoli).
\end{itemize}



\section{Correttura}
\begin{itemize}
\item Ogni esercizio va scritto su un foglio separato e numerato chiaramente. Non \`e consentito scrivere sul retro dei fogli.
\item Il candidato deve scrivere nome, cognome, numero dell'esercizio, numero di pagina e numero totale di pagine dell'esercizio. Per esempio: \textbf{Pietro Bernasconi, G1, pagina 2/4}.
\item Tutti i passaggi vanno argomentati rigorosamente. I teoremi presenti nei supporti didattici possono essere citati liberamente senza dimostrazione (ma la modalit\`a di applicazione deve essere trasparente). 
\item Dopo aver ricevuto i risultati, ogni candidato avr\`a la possibilit\`a di controllare la correzione effettuata e chiedere chiarimenti (entro dei limiti temporali che verranno indicati).
\end{itemize}
\end{document} 